\chapter{Introduction aux quantificateurs}
\label{chapter-quantifiers}

{\footnotesize\noindent
\review{Premier brouillon produit avec l'aide d'un modèle d'IA : ce
chapitre doit être relu, complété et vérifié par l'enseignant avant toute
diffusion.}\par}

\section{Les quantificateurs et les systèmes logiques}

\subsection{Les quantificateurs comme conjonctions et disjonctions en
famille}

\thingy
Jusqu'à présent, nous avons surtout considéré le \defin{calcul
propositionnel}, c'est-à-dire une logique dans laquelle les formules
représentent des affirmations et dans laquelle on peut les combiner par
les connecteurs propositionnels
\[
\Rightarrow,\qquad \land,\qquad \lor,\qquad \top,\qquad \bot.
\]
Ces connecteurs permettent déjà d'exprimer beaucoup de raisonnements.
Ils ne permettent toutefois pas d'exprimer directement des affirmations
du genre « tout entier possède telle propriété » ou « il existe un objet
possédant telle propriété ».

\thingy
Pour cela, on introduit deux notations logiques fondamentales : les
\defin{quantificateurs} universel et existentiel, notés respectivement
$\forall$ et $\exists$.  Ils s'appliquent à une formule $P(v)$ qui
dépend d'une variable libre $v$, supposée ici de type $I$.  Le
quantificateur lie alors cette variable et forme une nouvelle formule
\[
\forall (v:I).\,P(v)
\qquad\hbox{ou}\qquad
\exists (v:I).\,P(v).
\]
Lorsque le type des individus est clair ou laissé implicite, on écrit
souvent plus brièvement
\[
\forall v.\,P(v)
\qquad\hbox{et}\qquad
\exists v.\,P(v).
\]

\thingy
Intuitivement, on peut voir les quantificateurs comme des « $\land$ et
$\lor$ en famille ».  Plus précisément, la formule $\forall v.\,P(v)$,
parfois notée $\bigwedge_v P(v)$, joue le rôle d'une conjonction portant
sur toute une famille de formules.  De même, $\exists v.\,P(v)$, parfois
notée $\bigvee_v P(v)$, joue le rôle d'une disjonction portant sur une
famille.

Cette analogie est la même que celle qui relie les produits et les
sommes indexés aux produits et sommes binaires :
\[
\forall v.\,P(v)
\quad\hbox{est à}\quad
P\land Q
\quad\hbox{ce que}\quad
\prod_i p_i
\quad\hbox{est à}\quad
p\times q,
\]
et
\[
\exists v.\,P(v)
\quad\hbox{est à}\quad
P\lor Q
\quad\hbox{ce que}\quad
\sum_i p_i
\quad\hbox{est à}\quad
p+q.
\]

\thingy
Il faut cependant faire attention au fait que cette présentation reste
pour l'instant informelle.  Il existe de nombreux systèmes logiques,
qui diffèrent notamment par les objets sur lesquels on a le droit de
quantifier.  Autrement dit, il faudra préciser ce que peut être la
variable $v$, quel est son domaine ou son type $I$, et comment sont
formés les termes qui représentent des individus de ce type.

\subsection{Interprétation de Brouwer--Heyting--Kolmogorov}

\thingy
L'interprétation de Brouwer--Heyting--Kolmogorov, ou interprétation
\defin{BHK}, donne une lecture constructive des formules logiques : une
formule est comprise à travers ce qui constituerait un témoignage, ou
une preuve, de cette formule.  Pour les connecteurs propositionnels,
nous avons déjà vu les idées suivantes :
\begin{itemize}
\item un témoignage de $P\land Q$ est la donnée d'un témoignage de $P$
  et d'un témoignage de $Q$ ;
\item un témoignage de $P\lor Q$ est la donnée d'un témoignage de $P$
  ou d'un témoignage de $Q$, ainsi que l'indication du cas choisi ;
\item un témoignage de $P\Rightarrow Q$ est un moyen de transformer
  tout témoignage de $P$ en un témoignage de $Q$ ;
\item un témoignage de $\top$ est trivial ;
\item il n'existe pas de témoignage de $\bot$.
\end{itemize}

\thingy
Les quantificateurs prolongent naturellement cette lecture.  Un
témoignage de
\[
\forall v.\,P(v)
\]
est un moyen qui, à partir d'un individu quelconque $x$, fournit un
témoignage de $P(x)$.  En revanche, un témoignage de
\[
\exists v.\,P(v)
\]
est la donnée d'un certain individu $t_0$, accompagnée d'un témoignage
de $P(t_0)$.

Autrement dit, pour démontrer universellement une propriété, il faut
savoir traiter un individu arbitraire ; pour démontrer existentiellement
une propriété, il faut exhiber un individu qui convient.

\subsection{Introduction et élimination de $\forall$ et $\exists$}

\thingy
Pour introduire une quantification universelle
$\forall (v:I).\,Q(v)$, on choisit une variable $v:I$ arbitraire.  Il
faut que cette variable soit \defin{fraîche}, c'est-à-dire qu'elle
n'apparaisse librement dans aucune des hypothèses actuellement en cours.
Si l'on démontre alors $Q(v)$ pour ce $v$ arbitraire, on peut conclure
\[
\forall (v:I).\,Q(v).
\]

Dans une rédaction mathématique, le schéma correspondant est le
suivant : « soit $v$ arbitraire ; montrons $Q(v)$ ; donc
$\forall (v:I).\,Q(v)$ ».  La condition de fraîcheur est essentielle :
si $v$ dépendait d'une hypothèse particulière, il ne serait plus
véritablement arbitraire.

\thingy
Pour éliminer une quantification universelle, on part d'une preuve de
$\forall (v:I).\,Q(v)$ et l'on choisit un terme $t$ de type $I$.  On
obtient alors la formule $Q[v\backslash t]$, c'est-à-dire la formule
$Q$ dans laquelle les occurrences libres de $v$ ont été remplacées par
$t$ :
\[
\inferrule
  {\Gamma\vdash \forall (v:I).\,Q
   \\
   \Gamma\vdash t:I}
  {\Gamma\vdash Q[v\backslash t]}.
\]
En termes plus ordinaires, si une propriété est vraie pour tout
individu, elle est vraie en particulier pour l'individu représenté par
$t$.

\thingy
Pour introduire une quantification existentielle
$\exists (v:I).\,Q(v)$, il suffit au contraire d'exhiber un terme $t$
de type $I$, puis de démontrer $Q[v\backslash t]$.  On peut alors
conclure
\[
\inferrule
  {\Gamma\vdash t:I
   \\
   \Gamma\vdash Q[v\backslash t]}
  {\Gamma\vdash \exists (v:I).\,Q}.
\]
Le terme $t$ est le témoin de l'existence affirmée.

\thingy
Enfin, pour éliminer une quantification existentielle, on suppose que
l'on dispose de $\exists (v:I).\,P(v)$ et que l'on cherche à démontrer
une conclusion $Q$.  On peut alors considérer un individu $v:I$
arbitraire satisfaisant $P(v)$, à condition que $v$ soit frais : il ne
doit apparaître librement ni dans les hypothèses déjà présentes, ni
dans la conclusion $Q$.  Si, sous les hypothèses supplémentaires
$v:I$ et $P(v)$, on arrive à démontrer $Q$, alors on peut conclure $Q$
à partir de l'existence initiale :
\[
\inferrule
  {\Gamma\vdash \exists (v:I).\,P
   \\
   \Gamma,\,v:I,\,P\vdash Q}
  {\Gamma\vdash Q}.
\]

La condition portant sur $Q$ exprime que la conclusion ne doit pas
dépendre du témoin particulier choisi.  On raisonne bien à partir de
l'existence d'un témoin quelconque, et non à partir d'un témoin que
l'on aurait le droit de retenir en dehors de la démonstration.

\thingy
On peut résumer les règles de déduction naturelle ainsi obtenues par le
tableau suivant.  Le contexte $\Gamma$ peut contenir à la fois des
formules, qui jouent le rôle d'hypothèses logiques, et des variables
d'individus avec leur type, comme $v:I$.
\[
\begin{array}{c|c|c}
 & \text{Introduction} & \text{Élimination}\\\hline
\forall &
\inferrule{\Gamma,\,v:I\vdash Q}
          {\Gamma\vdash\forall(v:I).\,Q}
\quad(v\ \text{frais})
&
\inferrule{\Gamma\vdash\forall(v:I).\,Q\\\Gamma\vdash t:I}
          {\Gamma\vdash Q[v\backslash t]}
\\[1em]\hline
\exists &
\inferrule{\Gamma\vdash t:I\\\Gamma\vdash Q[v\backslash t]}
          {\Gamma\vdash\exists(v:I).\,Q}
&
\inferrule{\Gamma\vdash\exists(v:I).\,P\\
           \Gamma,\,v:I,\,P\vdash Q}
          {\Gamma\vdash Q}
\quad(v\ \text{frais})
\end{array}
\]

\subsection{Preuves, termes d'individus et notations de Curry--Howard}

\thingy
Dans l'interprétation de Curry--Howard, la conjonction $P\land Q$
correspond à un type produit $\sigma\times\tau$.  Une preuve de la
conjonction contient donc une preuve de chacune des deux composantes,
de même qu'une valeur d'un produit contient une valeur de chacun des
deux types.

Par analogie, la quantification universelle $\forall v.\,P(v)$
correspond à un \defin{type produit en famille}
\[
\prod_v \sigma(v).
\]
Un élément d'un tel type est une fonction qui, pour chaque $v$, fournit
une valeur de type $\sigma(v)$.  Une preuve de
$\forall(v:I).\,P(v)$ se note donc naturellement sous la forme
\[
\lambda(v:I).\,(\cdots),
\]
où le corps de l'abstraction est une preuve de $P(v)$.

\thingy
Avec les notations des $\lambda$-termes, la règle d'introduction de
$\forall$ devient
\[
\inferrule
  {\Gamma,\,v:I\vdash s:Q}
  {\Gamma\vdash\lambda(v:I).\,s
   \;:\;\forall(v:I).\,Q}.
\]
La variable $v$ peut apparaître dans $Q$ et dans le terme $s$, mais elle
ne doit pas apparaître dans le contexte $\Gamma$.

L'élimination de $\forall$ est alors une application ordinaire :
\[
\inferrule
  {\Gamma\vdash f:\forall(v:I).\,Q
   \\
   \Gamma\vdash t:I}
  {\Gamma\vdash ft:Q[v\backslash t]}.
\]
Une preuve de $\forall(v:I).\,Q(v)$ se comporte donc comme une fonction
qui reçoit un individu $t$ de type $I$ et renvoie une preuve de $Q(t)$.

\thingy
De même, la disjonction $P\lor Q$ correspond à un type somme
$\sigma+\tau$.  La quantification existentielle $\exists v.\,P(v)$
correspond alors à un \defin{type somme en famille}
\[
\sum_v \sigma(v).
\]
Une valeur de ce type est la donnée d'un indice $t_0$, accompagnée
d'une valeur de type $\sigma(t_0)$.  Une preuve de
$\exists(v:I).\,P(v)$ prend ainsi la forme d'une paire
\[
\langle t_0,\cdots\rangle.
\]

\thingy
La règle d'introduction de $\exists$ peut se noter
\[
\inferrule
  {\Gamma\vdash t:I
   \\
   \Gamma\vdash z:Q[v\backslash t]}
  {\Gamma\vdash\langle t,z\rangle
   \;:\;\exists(v:I).\,Q}.
\]
Le premier composant est le témoin, et le second est la preuve que ce
témoin satisfait la propriété demandée.

Pour éliminer une preuve existentielle, on utilise un filtrage de motif
analogue à celui employé pour les types sommes :
\[
\inferrule
  {\Gamma\vdash z:\exists(v:I).\,P
   \\
   \Gamma,\,v:I,\,h:P\vdash s:Q}
  {\Gamma\vdash
   (\texttt{match~}z\texttt{~with~}\langle v,h\rangle\mapsto s)
   \;:\;Q}.
\]
Ici encore, $v$ peut apparaître dans $s$, mais ne doit apparaître ni
dans $\Gamma$ ni dans $Q$.

\thingy
Cette correspondance entre quantificateurs et produits ou sommes en
famille a toutefois une limite importante.  Il n'est pas toujours
possible, dans un système logique donné, d'extraire le témoin $t_0$
d'une preuve de $\exists v.\,P(v)$ afin de l'utiliser comme une donnée
ordinaire.  Dans certains systèmes, le témoin ne peut être exploité
qu'à l'intérieur d'une preuve logique.

\review{La portée exacte de cette restriction dépend du système de types
considéré.  Elle devra être vérifiée lors de l'harmonisation du présent
chapitre avec les chapitres consacrés à Curry--Howard et aux systèmes de
types dépendants.}

\subsection{Ce que l'on peut quantifier : systèmes de types et
$\lambda$-cube}

\thingy
Les règles précédentes ne suffisent pas encore à définir complètement
un système logique.  Il faut aussi préciser la formation des termes
d'individus, ainsi que la façon dont le monde des individus et le monde
des propositions interagissent.

On peut distinguer, selon les systèmes, deux mondes plus ou moins
séparés :
\begin{itemize}
\item le monde des termes, ou individus, et de leurs types ;
\item le monde logique, constitué des propositions et de leurs preuves.
\end{itemize}
Les règles de déduction naturelle présentées ci-dessus construisent des
démonstrations et appartiennent donc au monde logique.  Les prémisses de
la forme $\Gamma\vdash t:I$, qui assurent qu'un terme $t$ est bien un
individu de type $I$, relèvent du monde des termes.

\thingy
Dans Coq/Rocq, ces deux mondes sont séparés mais parallèles : les
propositions appartiennent à $\mathit{Prop}$, tandis que les types
d'individus appartiennent à $\mathit{Type}$.  D'autres systèmes adoptent
des séparations différentes, ou mélangent davantage ces deux niveaux.

Les questions suivantes permettent notamment de distinguer les
systèmes :
\begin{itemize}
\item peut-on quantifier sur les propositions, comme dans
  \[
  \forall(A:*).\,\bigl(A\Rightarrow A\bigr),
  \]
  avec la preuve
  $\lambda(A:*).\,\lambda(h:A).\,h$ ?
\item peut-on abstraire sur une proposition, et former par exemple
  \[
  \lambda(A:*).\,A
  \]
  de type $*\rightarrow *$ ?
\item peut-on quantifier sur des individus $x:I$ dans une famille de
  propositions $A:I\rightarrow *$, comme dans
  \[
  \forall(x:I).\,A(x)?
  \]
\end{itemize}

\thingy
Les réponses possibles à ces questions conduisent notamment aux huit
systèmes du \defin{$\lambda$-cube} de Barendregt.  Le but n'est pas ici
de les étudier en détail, mais de retenir que les règles de quantification
dépendent étroitement de ce que le système autorise comme objets, types,
propositions et abstractions.

\subsection{Quantification existentielle, types sommes et
imprédicativité}

\thingy
Considérons l'énoncé suivant, pour deux types $U$ et $V$ :
\[
\bigl(\forall(x:U).\,\exists(y:V).\,P(x,y)\bigr)
\Rightarrow
\bigl(\exists(f:U\Rightarrow V).\,\forall(x:U).\,P(x,f(x))\bigr).
\]
Cet énoncé est un analogue, dans le langage des types, de l'\defin{axiome
du choix}.  Il affirme que, si pour chaque $x:U$ on dispose d'un $y:V$
tel que $P(x,y)$, alors on peut réunir les témoins ainsi obtenus dans une
fonction $f:U\Rightarrow V$.

\thingy
Si $\forall$ et $\exists$ sont interprétés respectivement comme un
produit et une somme en famille, cet énoncé paraît naturel.  Une preuve
de
\[
\forall(x:U).\,\exists(y:V).\,P(x,y)
\]
est alors une fonction qui, à chaque $x$, associe une paire constituée
d'un $y$ et d'une preuve de $P(x,y)$.  On peut donc définir la fonction
qui extrait le premier composant de cette paire, et conserver dans une
seconde preuve le second composant.

En revanche, si $\exists$ est seulement un quantificateur logique, le
témoin fourni par une preuve existentielle ne peut pas nécessairement
être extrait comme une valeur utilisable dans le monde des programmes.
L'énoncé ci-dessus n'est alors pas nécessairement démontrable.

\thingy
Cette différence est particulièrement visible entre des systèmes tels
que Coq, où l'énoncé précédent n'est pas en général démontrable lorsque
$P:U\times V\to\mathit{Prop}$, et des systèmes à la Martin-Löf tels
qu'Agda, dans lesquels les preuves existentielles sont des paires
effectives et permettent cette construction.

\review{Cette comparaison entre Coq/Rocq et Agda reprend l'idée des
diapositives, mais devra être vérifiée avec soin au regard des univers,
des restrictions d'élimination de \texttt{Prop} et des bibliothèques
effectivement utilisées.}

\thingy
On appelle \defin{imprédicativité} la possibilité de définir une
proposition ou un type en quantifiant sur toutes les propositions ou
tous les types, y compris celui que l'on est en train de définir.  Il
s'agit donc d'une certaine forme de circularité.

Par exemple, en notant $*$ un « type des types » imprédicatif, on peut
considérer les expressions suivantes :
\[
A \;\cong\; \forall(Z:*).\,\bigl(Z\Rightarrow A\bigr),
\]
\[
A \;\cong\;
\forall(Z:*).\,\bigl((A\Rightarrow Z)\Rightarrow Z\bigr),
\]
\[
\bot\;\cong\;\forall(Z:*).\,Z,
\]
\[
A\land B\;\cong\;
\forall(Z:*).\,\bigl((A\Rightarrow B\Rightarrow Z)\Rightarrow Z\bigr),
\]
et
\[
A\lor B\;\cong\;
\forall(Z:*).\,\bigl((A\Rightarrow Z)\Rightarrow
(B\Rightarrow Z)\Rightarrow Z\bigr).
\]
Ces équivalences sont à comprendre ici comme des indications
informelles de codages possibles grâce à la quantification
imprédicative.

\thingy
L'imprédicativité est utile pour définir certaines constructions sur les
types.  Elle doit néanmoins être maniée avec prudence : sous des formes
trop générales, elle peut mener à des incohérences logiques, comme le
montre le paradoxe de Girard.

\review{Les codages imprédicatifs ci-dessus sont seulement esquissés dans
les diapositives.  Une passe ultérieure devra décider s'il convient de
donner les termes précis des deux sens de chaque isomorphisme, ou de
conserver cette présentation informelle.}

\section{Logique du premier ordre}

\subsection{Individus, prédicats et syntaxe des formules}

\thingy
La \defin{logique du premier ordre}, aussi appelée \defin{calcul des
prédicats} du premier ordre, est le système logique le plus simple qui
ajoute les quantificateurs au calcul propositionnel.  Les objets sur
lesquels il est permis de quantifier sont appelés des \defin{individus}.

Du point de vue des systèmes de types, cette logique est assez
rudimentaire : les individus forment un unique type implicite $I$,
spécialement réservé à ce rôle.  La logique pure du premier ordre ne
permet pas directement de construire des couples, des fonctions ou
d'autres objets à partir de ces individus.  Cette limitation est
compensée, dans les usages mathématiques ordinaires, par la théorie que
l'on exprime dans cette logique.

\thingy
La logique du premier ordre possède une importance mathématique
considérable.  La position classique est que les mathématiques peuvent
être formalisées dans la théorie des ensembles de Zermelo--Fraenkel avec
l'axiome du choix, notée $\mathsf{ZFC}$, elle-même formulée en logique
du premier ordre.  Les objets mathématiques élaborés, tels que couples,
fonctions ou structures, sont alors construits à l'intérieur de la
théorie des ensembles plutôt que fournis directement par la logique.

\thingy
Dans la logique du premier ordre pure, on dispose de deux sortes de
variables :
\begin{itemize}
\item des variables d'individus, notées par exemple $x$, $y$, $z$,
  en nombre illimité ;
\item pour chaque entier naturel $n$, des variables de prédicats
  $n$-aires, aussi appelées relations $n$-aires entre individus, notées
  par exemple $A^{(n)}$, $B^{(n)}$, $C^{(n)}$.
\end{itemize}
L'arité est généralement omise lorsqu'elle se lit immédiatement sur la
formule.

\thingy
Les formules sont définies inductivement.  Une formule est :
\begin{itemize}
\item une application $A^{(n)}(x_1,\ldots,x_n)$ d'une variable de
  prédicat $n$-aire à $n$ variables d'individus ;
\item une combinaison de deux formules $P$ et $Q$ par l'un des
  connecteurs $(P\Rightarrow Q)$, $(P\land Q)$ ou $(P\lor Q)$, ou bien
  l'une des constantes $\top$ et $\bot$ ;
\item une quantification $\forall x.\,P$ ou $\exists x.\,P$, qui lie la
  variable d'individu $x$ dans $P$.
\end{itemize}
Le type $I$ des individus reste implicite : les notations
$\forall x.\,P$ et $\exists x.\,P$ abrègent donc
$\forall(x:I).\,P$ et $\exists(x:I).\,P$.

\thingy
Il ne faut pas confondre les variables de prédicats avec les variables
d'individus.  Dans la logique du premier ordre, on peut quantifier sur
les individus, mais non sur les prédicats eux-mêmes : c'est précisément
le sens de l'expression « premier ordre ».

\subsection{Règles de déduction naturelle et conditions de fraîcheur}

\thingy
Les règles logiques de la logique du premier ordre reprennent les règles
du calcul propositionnel, auxquelles s'ajoutent les règles
d'introduction et d'élimination de $\forall$ et $\exists$.  Elles sont
les mêmes que celles données en
\ref{section-quantifiers-and-logical-systems}, mais la variable
quantifiée est désormais une variable d'individu $x$ du type implicite
$I$.

\review{Le label
\texttt{section-quantifiers-and-logical-systems} doit être ajouté à la
première section si l'on souhaite conserver ce renvoi.}

\thingy
Un point essentiel est que les prémisses de formation des termes,
comme
\[
\Gamma\vdash t:I,
\]
ne doivent pas être oubliées.  Elles assurent que le terme substitué
dans une formule universelle ou existentielle désigne effectivement un
individu.  Dans les arbres de preuve, ces prémisses sont parfois
omises pour alléger la présentation, mais elles restent nécessaires.

\thingy
Dans la variante la plus simple de la logique du premier ordre, les
seuls termes d'individus sont les variables d'individus.  Ainsi, la
seule façon d'obtenir une dérivation de $t:I$ est essentiellement
d'avoir déclaré $t:I$ dans le contexte.  La règle d'axiome doit donc
être comprise sous la forme
\[
\inferrule{}{\,\Gamma,\,x:I\vdash x:I}
\]
et, pour une formule atomique,
\[
\inferrule
  {\Gamma\vdash x_1:I\\\cdots\\\Gamma\vdash x_n:I}
  {\Gamma,\,A(x_1,\ldots,x_n)\vdash A(x_1,\ldots,x_n)}.
\]

\thingy
On ne suppose pas que le domaine des individus soit habité.  En
particulier, la formule
\[
\exists x.\,\top
\]
n'est pas démontrable sans hypothèse supplémentaire.  Pour la même
raison, on ne peut pas démontrer en général
\[
\bigl(\forall x.\,A(x)\bigr)\Rightarrow\bigl(\exists x.\,A(x)\bigr).
\]
Une preuve universelle ne fournit pas à elle seule un individu auquel
l'appliquer.

\subsection{Exemples de démonstrations et notation des preuves}

\thingy
Les formules propositionnelles sont des cas particuliers de formules du
premier ordre.  Il suffit, pour cela, de considérer les variables
propositionnelles comme des prédicats $0$-aires, aussi dits
\defin{nullaires}.  Ainsi,
\[
A\land B\Rightarrow B\land A
\]
est à la fois une formule propositionnelle et une formule du premier
ordre.

\thingy
Voici quelques exemples de formules du premier ordre :
\begin{itemize}
\item
  \[
  (\forall x.\,A(x))\land(\exists x.\,\top)
  \Rightarrow(\exists x.\,A(x));
  \]
\item
  \[
  (\forall x.\,\neg A(x))\Leftrightarrow(\neg\exists x.\,A(x));
  \]
\item
  \[
  (\exists x.\,\neg A(x))\Rightarrow(\neg\forall x.\,A(x));
  \]
\item
  \[
  (\exists x.\,A)\Leftrightarrow(\exists x.\,\top)\land A,
  \]
  où $A$ est un prédicat nullaire ;
\item
  \[
  (\forall x.\,A)\Leftrightarrow((\exists x.\,\top)\Rightarrow A),
  \]
  avec la même convention ;
\item
  \[
  (\exists x.\,\forall y.\,B(x,y))
  \Rightarrow
  (\forall y.\,\exists x.\,B(x,y)).
  \]
\end{itemize}

\thingy
Dans ces notes, nous suivons la convention selon laquelle les
quantificateurs $\forall$ et $\exists$ ont une priorité plus faible que
les connecteurs $\Rightarrow$, $\lor$, $\land$ et $\neg$.  Cette
convention n'est pas universelle : lorsque la moindre ambiguïté est
possible, il est préférable d'ajouter des parenthèses.

\thingy
Démontrons l'implication
\[
(\exists x.\,\forall y.\,B(x,y))
\Rightarrow
(\forall y.\,\exists x.\,B(x,y)).
\]
Supposons donc $\exists x.\,\forall y.\,B(x,y)$.  Soit $y'$ un individu
arbitraire.  Par élimination de l'existentielle, considérons un
individu $x$ tel que $\forall y.\,B(x,y)$.  Par élimination de
l'universelle, on obtient $B(x,y')$.  En introduisant l'existentielle
avec le témoin $x$, on en déduit $\exists x'.\,B(x',y')$.

Ce raisonnement ne dépend pas du choix de $y'$ ; on peut donc introduire
$\forall y'$.  De même, le témoin $x$ utilisé dans le raisonnement
intermédiaire était obtenu à partir de l'hypothèse existentielle ; on
peut donc conclure sans le conserver hors de cette élimination.  Enfin,
on décharge l'hypothèse initiale pour introduire l'implication.

\thingy
Avec les notations de Curry--Howard, la preuve précédente s'écrit
\[
\lambda(u:\exists x.\,\forall y.\,B(x,y)).\;
\lambda(y':I).\;
(\texttt{match~}u\texttt{~with~}\langle x,v\rangle\mapsto
\langle x,vy'\rangle).
\]
Le terme $u$ contient une paire existentielle : le filtrage de motif en
extrait un témoin $x$ et une preuve $v$ de $\forall y.\,B(x,y)$.  En
appliquant $v$ à $y'$, on obtient $vy':B(x,y')$, puis la paire
$\langle x,vy'\rangle$ constitue une preuve de
$\exists x'.\,B(x',y')$.

\subsection{Le monde des individus et le monde logique}

\thingy
La logique du premier ordre maintient une séparation particulièrement
nette entre le type implicite $I$ des individus et les relations
$A,B,C,\ldots$ qui servent à former des propositions.  Le type $I$ est
spécial : on ne l'écrit généralement pas, et la logique pure ne permet
pas de former directement des types tels que $I\times I$ ou $I\to I$.

\thingy
Les individus ne se mélangent donc pas directement aux preuves.  Les
variables $x$, $y$, $z$ appartiennent au monde des individus, tandis
que les formules et leurs démonstrations appartiennent au monde
logique.  Les prédicats permettent de parler logiquement des individus,
mais ils ne transforment pas pour autant les individus en propositions
ni les preuves en individus.

Cette organisation explique à la fois la simplicité de la logique du
premier ordre et ses limitations du point de vue du typ
